% book11.tex --- Book XI of Euclid's Elements: Solid Geometry.
%
% All 39 propositions encoded.  Book XI extends the plane geometry of
% Books I-IV to three dimensions: lines and planes in space, parallel
% planes, dihedral angles, parallelepipeds, and the foundational lemmas
% for the Platonic solids of Book XIII.
%
% Wording follows Heath (1908).

\section{Book XI --- Solid Geometry}
\label{sec:book-XI}

\begin{claim}[Proposition XI.1: A straight line cannot have part in a plane and part not]
\label{prop:XI.1}
A part of a straight line cannot be in the plane of reference and a
part in a plane more elevated.
\end{claim}
\begin{evidence}[Proof of XI.1]
\label{ev:XI.1}
If a straight line $AB$ had part in one plane and continued part in
another, then through $B$ there would be two distinct straight lines
from $A$ (one in each plane), contradicting Postulate 1 (uniqueness
of the straight line through two points).
\dependson{XI.1}{post:1}
\dependson{XI.1}{def:I.4}
\end{evidence}

\begin{claim}[Proposition XI.2: Intersecting lines determine a plane]
\label{prop:XI.2}
If two straight lines cut one another, they are in one plane, and
every triangle is in one plane.
\end{claim}
\begin{evidence}[Proof of XI.2]
\label{ev:XI.2}
Two intersecting lines pick out three non-collinear points (one at
the intersection, one on each line); through these three points
passes exactly one plane (the analogue in 3D of Postulate 1).
\dependson{XI.2}{XI.1}
\dependson{XI.2}{def:I.7}
\end{evidence}

\begin{claim}[Proposition XI.3: Intersection of two planes is a straight line]
\label{prop:XI.3}
If two planes cut one another, their common section is a straight
line.
\end{claim}
\begin{evidence}[Proof of XI.3]
\label{ev:XI.3}
Take two points $A$, $B$ in the common section.  Draw the straight
line $AB$ in each plane; by XI.1 each segment of $AB$ lies in its
plane, and uniqueness of the line forces both segments to coincide.
\dependson{XI.3}{XI.1}
\dependson{XI.3}{post:1}
\end{evidence}

\begin{claim}[Proposition XI.4: Line perpendicular to two intersecting lines is perpendicular to their plane]
\label{prop:XI.4}
If a straight line be set up at right angles to two straight lines
which cut one another, at their common point of section, it will
also be at right angles to the plane through them.
\end{claim}
\begin{evidence}[Proof of XI.4]
\label{ev:XI.4}
Take any other straight line $\ell$ through the foot in the plane;
$\ell$ can be expressed as a sum of perpendicular components on the
two given lines (by I.46-style decomposition), and the perpendicular
to both is perpendicular to the sum by I.4 applied to the right
triangles formed.
\dependson{XI.4}{I.4}
\dependson{XI.4}{I.8}
\dependson{XI.4}{def:XI.3}
\end{evidence}

\begin{claim}[Proposition XI.5: Three lines through a point each perpendicular to a fourth lie in a plane]
\label{prop:XI.5}
If a straight line be set up at right angles to three straight lines
which meet one another, at their common point of section, the three
straight lines are in one plane.
\end{claim}
\begin{evidence}[Proof of XI.5]
\label{ev:XI.5}
Two of the three meeting lines determine a plane (XI.2); if the
third were out of that plane, the perpendicular relation combined
with XI.4 would force two distinct planes through the same set of
perpendicular lines, contradiction.
\dependson{XI.5}{XI.2}
\dependson{XI.5}{XI.4}
\end{evidence}

\begin{claim}[Proposition XI.6: Two lines perpendicular to the same plane are parallel]
\label{prop:XI.6}
If two straight lines be at right angles to the same plane, the
straight lines will be parallel.
\end{claim}
\begin{evidence}[Proof of XI.6]
\label{ev:XI.6}
Suppose the two perpendiculars $AB$, $CD$ met or were skew.  Drop
the segment $BD$ in the plane; the angles at $B$ and $D$ are right.
In the plane through $AB$ and $BD$, by I.28 the line $CD$ would
need to be parallel to $AB$, fixing the plane through $AB$, $CD$.
Then within that plane the two right angles force $AB \parallel CD$.
\dependson{XI.6}{I.28}
\dependson{XI.6}{XI.2}
\dependson{XI.6}{XI.3}
\dependson{XI.6}{XI.4}
\end{evidence}

\begin{claim}[Proposition XI.7: A line in a plane parallel to a parallel line in another plane lies in the connecting plane]
\label{prop:XI.7}
If two straight lines be parallel, and points be taken at random on
each of them, the straight line joining the points is in the same
plane with the parallels.
\end{claim}
\begin{evidence}[Proof of XI.7]
\label{ev:XI.7}
Two parallel lines determine a plane (XI.2 extended).  Any joining
segment between points on the two parallels lies in this plane by
XI.1.
\dependson{XI.7}{XI.1}
\dependson{XI.7}{XI.2}
\end{evidence}

\begin{claim}[Proposition XI.8: Two parallel lines, one perpendicular to a plane, the other also perpendicular]
\label{prop:XI.8}
If two straight lines be parallel, and one of them be at right angles
to any plane, the remaining one will also be at right angles to the
same plane.
\end{claim}
\begin{evidence}[Proof of XI.8]
\label{ev:XI.8}
By XI.7 the connecting segment lies in the plane of the parallels;
combine XI.4 (perpendicularity transferred along parallel lines via
shared plane structure) with I.29 (alternate angles) to obtain
perpendicularity of the second line.
\dependson{XI.8}{I.29}
\dependson{XI.8}{XI.4}
\dependson{XI.8}{XI.6}
\dependson{XI.8}{XI.7}
\end{evidence}

\begin{claim}[Proposition XI.9: Lines parallel to the same line are parallel]
\label{prop:XI.9}
Straight lines which are parallel to the same straight line and are
not in the same plane with it are also parallel to one another.
\end{claim}
\begin{evidence}[Proof of XI.9]
\label{ev:XI.9}
Construct in each plane perpendiculars from a common point to the
shared straight line; the perpendiculars are equal in length, and
by I.33 the resulting transversal is parallel to both targets.
\dependson{XI.9}{I.33}
\dependson{XI.9}{XI.6}
\dependson{XI.9}{XI.8}
\end{evidence}

\begin{claim}[Proposition XI.10: Two angles with parallel sides are equal in space]
\label{prop:XI.10}
If two straight lines meeting one another be parallel to two straight
lines meeting one another, not in the same plane, they will contain
equal angles.
\end{claim}
\begin{evidence}[Proof of XI.10]
\label{ev:XI.10}
Construct the parallelogram joining corresponding points; by I.33
opposite sides are equal, and by I.8 (SSS) the two triangles formed
at the vertex angles are congruent.
\dependson{XI.10}{I.8}
\dependson{XI.10}{I.33}
\dependson{XI.10}{XI.9}
\end{evidence}

\begin{claim}[Proposition XI.11: Drop a perpendicular from an external point to a plane]
\label{prop:XI.11}
From a given elevated point to draw a straight line perpendicular to
a given plane.
\end{claim}
\begin{evidence}[Proof of XI.11]
\label{ev:XI.11}
Drop a chord through the point parallel to the plane; drop a
perpendicular from the chord to its foot in the plane; the
constructed line, being perpendicular to two intersecting lines at
the foot, is perpendicular to the plane (XI.4).
\dependson{XI.11}{I.11}
\dependson{XI.11}{I.12}
\dependson{XI.11}{XI.4}
\end{evidence}

\begin{claim}[Proposition XI.12: Erect a perpendicular to a plane at a given point in the plane]
\label{prop:XI.12}
To set up a straight line at right angles to a given plane from a
given point in it.
\end{claim}
\begin{evidence}[Proof of XI.12]
\label{ev:XI.12}
Take an external point above the plane; drop a perpendicular from it
to the plane via XI.11; the foot may be made to coincide with the
given point by an additional parallel translation (using XI.8).
\dependson{XI.12}{XI.8}
\dependson{XI.12}{XI.11}
\end{evidence}

\begin{claim}[Proposition XI.13: A unique perpendicular to a plane from a given point]
\label{prop:XI.13}
From the same point two straight lines cannot be set up at right
angles to the same plane on the same side.
\end{claim}
\begin{evidence}[Proof of XI.13]
\label{ev:XI.13}
Two such perpendiculars would meet two lines in the plane at the
same right angles; by XI.4 / XI.6 the two perpendiculars would have
to be parallel; but parallels do not meet -- contradicting their
common origin.
\dependson{XI.13}{XI.4}
\dependson{XI.13}{XI.6}
\end{evidence}

\begin{claim}[Proposition XI.14: Planes perpendicular to the same line are parallel]
\label{prop:XI.14}
Planes to which the same straight line is at right angles will be
parallel.
\end{claim}
\begin{evidence}[Proof of XI.14]
\label{ev:XI.14}
If the two planes met, the line of intersection (XI.3) would meet
the common perpendicular at right angles in two distinct places --
contradicting XI.13.
\dependson{XI.14}{XI.3}
\dependson{XI.14}{XI.13}
\dependson{XI.14}{def:XI.8}
\end{evidence}

\begin{claim}[Proposition XI.15: Two intersecting line-pairs parallel to two other line-pairs give parallel planes]
\label{prop:XI.15}
If two straight lines meeting one another be parallel to two straight
lines meeting one another, not being in the same plane, the planes
through them are parallel.
\end{claim}
\begin{evidence}[Proof of XI.15]
\label{ev:XI.15}
By XI.10 the angles are equal; drop a common perpendicular line; by
XI.14 the two planes share a common perpendicular and are parallel.
\dependson{XI.15}{XI.9}
\dependson{XI.15}{XI.10}
\dependson{XI.15}{XI.14}
\end{evidence}

\begin{claim}[Proposition XI.16: A plane cuts parallel planes in parallel lines]
\label{prop:XI.16}
If two parallel planes be cut by any plane, their common sections
are parallel.
\end{claim}
\begin{evidence}[Proof of XI.16]
\label{ev:XI.16}
The two intersection lines lie in the cutting plane, and if they
met, the meeting point would belong to both parallel planes -- which
is impossible.
\dependson{XI.16}{XI.3}
\dependson{XI.16}{def:XI.8}
\end{evidence}

\begin{claim}[Proposition XI.17: Two parallel planes cut a transversal proportionally]
\label{prop:XI.17}
If two straight lines be cut by parallel planes, they will be cut in
the same ratios.
\end{claim}
\begin{evidence}[Proof of XI.17]
\label{ev:XI.17}
Draw a parallelogram structure between the parallel planes; apply
VI.2 (basic proportionality) in each pair of cross-sectional lines.
\dependson{XI.17}{VI.2}
\dependson{XI.17}{XI.16}
\end{evidence}

\begin{claim}[Proposition XI.18: A line perpendicular to a plane makes the plane through it perpendicular to that plane]
\label{prop:XI.18}
If a straight line be at right angles to any plane, all the planes
through it will also be at right angles to the same plane.
\end{claim}
\begin{evidence}[Proof of XI.18]
\label{ev:XI.18}
Any plane through the perpendicular contains the perpendicular line;
by Definition XI.4 (perpendicular planes) the containing plane is
perpendicular to the original plane.
\dependson{XI.18}{XI.4}
\dependson{XI.18}{def:XI.4}
\end{evidence}

\begin{claim}[Proposition XI.19: Two perpendicular planes intersect in a line perpendicular to the base]
\label{prop:XI.19}
If two planes which cut one another be at right angles to any plane,
their common section will also be at right angles to the same plane.
\end{claim}
\begin{evidence}[Proof of XI.19]
\label{ev:XI.19}
The intersection line lies in both planes; by Definition XI.4 the
perpendiculars from any point of the intersection within each plane
are perpendicular to the base plane; XI.13 then forces the
intersection line itself to be perpendicular.
\dependson{XI.19}{XI.13}
\dependson{XI.19}{XI.18}
\dependson{XI.19}{def:XI.4}
\end{evidence}

\begin{claim}[Proposition XI.20: Solid angle inequality (triangle inequality for face-angles)]
\label{prop:XI.20}
If a solid angle be contained by three plane angles, any two, taken
together in any manner, are greater than the remaining one.
\end{claim}
\begin{evidence}[Proof of XI.20]
\label{ev:XI.20}
Suppose the largest face-angle is $\angle BAC$.  Within $\angle BAC$
construct $\angle BAD$ equal to $\angle BAE$ (one of the other
face-angles).  By I.4 / I.24, the corresponding chord arcs in space
give the desired strict triangle-style inequality among the
face-angles.
\dependson{XI.20}{I.4}
\dependson{XI.20}{I.20}
\dependson{XI.20}{I.24}
\end{evidence}

\begin{claim}[Proposition XI.21: Sum of face-angles at a solid angle is less than four right angles]
\label{prop:XI.21}
Any solid angle is contained by plane angles less than four right
angles.
\end{claim}
\begin{evidence}[Proof of XI.21]
\label{ev:XI.21}
Cut a small polygon by a plane near the apex; the sum of the
exterior angles of this polygon is less than $4 \cdot 90^\circ$ (by
I.32 / I.34 applied to the polygon).  The interior face-angles at
the apex are the supplements of these exterior angles, so their sum
falls strictly short of $4 \cdot 90^\circ$.
\dependson{XI.21}{I.32}
\dependson{XI.21}{XI.20}
\end{evidence}

\begin{claim}[Proposition XI.22: Three face-angles whose sum is less than four right angles can form a solid angle]
\label{prop:XI.22}
If there be three plane angles of which two, taken together in any
manner, are greater than the remaining one, and they are contained
by equal straight lines, it is possible to construct a triangle out
of the straight lines joining the extremities of the equal straight
lines.
\end{claim}
\begin{evidence}[Proof of XI.22]
\label{ev:XI.22}
The plane-angle inequality (XI.20) is precisely the triangle
inequality for the joining chords; by I.22 (construction of a
triangle on three given segments) the resulting triangle exists.
\dependson{XI.22}{I.22}
\dependson{XI.22}{XI.20}
\end{evidence}

\begin{claim}[Proposition XI.23: Construct a solid angle from three given face-angles]
\label{prop:XI.23}
To construct a solid angle out of three plane angles, two of which,
taken together in any manner, are greater than the remaining one;
thus the sum of the three angles must be less than four right angles.
\end{claim}
\begin{evidence}[Proof of XI.23]
\label{ev:XI.23}
Apply XI.22 to obtain the triangle of chords; mount that triangle so
that the three face-angles meet at a common apex; the bound of XI.21
ensures consistency.
\dependson{XI.23}{XI.20}
\dependson{XI.23}{XI.21}
\dependson{XI.23}{XI.22}
\end{evidence}

\begin{claim}[Proposition XI.24: Parallelepiped has parallelogram faces]
\label{prop:XI.24}
If a solid be contained by parallel planes, the opposite planes in
it are equal and similar parallelograms.
\end{claim}
\begin{evidence}[Proof of XI.24]
\label{ev:XI.24}
Opposite faces share parallel sides (by XI.16) and equal angles (by
XI.10), so they are congruent parallelograms by I.33 / I.34.
\dependson{XI.24}{I.33}
\dependson{XI.24}{I.34}
\dependson{XI.24}{XI.10}
\dependson{XI.24}{XI.16}
\end{evidence}

\begin{claim}[Proposition XI.25: Parallelepiped cut by a plane parallel to a face is divided proportionally]
\label{prop:XI.25}
If a parallelepipedal solid be cut by a plane parallel to opposite
planes, then, as the base is to the base, so will the solid be to
the solid.
\end{claim}
\begin{evidence}[Proof of XI.25]
\label{ev:XI.25}
Apply XI.17 to the side faces and VI.1 to the parallel base-pairs;
the volume is proportional to one varying side at constant
cross-section.
\dependson{XI.25}{VI.1}
\dependson{XI.25}{XI.17}
\dependson{XI.25}{XI.24}
\end{evidence}

\begin{claim}[Proposition XI.26: Construct a solid angle equal to a given solid angle]
\label{prop:XI.26}
At a given point on a given straight line to construct a solid angle
equal to a given solid angle contained by three plane angles.
\end{claim}
\begin{evidence}[Proof of XI.26]
\label{ev:XI.26}
Reproduce each face-angle by I.23 in the appropriate planes; by XI.23
the resulting figure determines a solid angle congruent to the given
one.
\dependson{XI.26}{I.23}
\dependson{XI.26}{XI.10}
\dependson{XI.26}{XI.23}
\end{evidence}

\begin{claim}[Proposition XI.27: Construct a parallelepiped similar to a given parallelepiped on a given edge]
\label{prop:XI.27}
On a given straight line to construct a parallelepipedal solid
similar and similarly situated to a given parallelepipedal solid.
\end{claim}
\begin{evidence}[Proof of XI.27]
\label{ev:XI.27}
Apply the face-angle construction of XI.26 at each vertex; by VI.18
the resulting faces are similar to the corresponding faces of the
given solid.
\dependson{XI.27}{VI.18}
\dependson{XI.27}{XI.24}
\dependson{XI.27}{XI.26}
\end{evidence}

\begin{claim}[Proposition XI.28: Parallelepiped bisected by its diagonal plane]
\label{prop:XI.28}
If a parallelepipedal solid be cut by a plane through the diagonals
of the opposite planes, the solid will be bisected by the plane.
\end{claim}
\begin{evidence}[Proof of XI.28]
\label{ev:XI.28}
The two pieces are mirror-image prisms with congruent base-triangles
(by I.34); by XI.24 their volumes are equal.
\dependson{XI.28}{I.34}
\dependson{XI.28}{XI.24}
\dependson{XI.28}{XI.25}
\end{evidence}

\begin{claim}[Proposition XI.29: Parallelepipeds on equal bases with same height are equal]
\label{prop:XI.29}
Parallelepipedal solids which are on the same base and of the same
height, and in which the extremities of the sides which stand up are
on the same straight lines, are equal to one another.
\end{claim}
\begin{evidence}[Proof of XI.29]
\label{ev:XI.29}
The two parallelepipeds can be decomposed into congruent prisms via
XI.28 and I.34 applied repeatedly; the equality is the 3D analogue
of I.35 (parallelograms on the same base between the same parallels).
\dependson{XI.29}{I.35}
\dependson{XI.29}{XI.24}
\dependson{XI.29}{XI.28}
\end{evidence}

\begin{claim}[Proposition XI.30: Parallelepipeds on equal bases with same height but different oblique placement are equal]
\label{prop:XI.30}
Parallelepipedal solids which are on the same base and of the same
height, and in which the extremities of the sides which stand up are
not on the same straight lines, are equal to one another.
\end{claim}
\begin{evidence}[Proof of XI.30]
\label{ev:XI.30}
Variant of XI.29 with the oblique sides at different angles; the
decomposition argument still applies after a shear.
\dependson{XI.30}{XI.29}
\end{evidence}

\begin{claim}[Proposition XI.31: Parallelepipeds on equal bases are in the ratio of their heights]
\label{prop:XI.31}
Parallelepipedal solids which are on equal bases and of the same
height are equal to one another.
\end{claim}
\begin{evidence}[Proof of XI.31]
\label{ev:XI.31}
By XI.25 the volume is proportional to the base at fixed height; if
the bases are equal, the volumes are equal.
\dependson{XI.31}{XI.25}
\dependson{XI.31}{XI.29}
\dependson{XI.31}{XI.30}
\end{evidence}

\begin{claim}[Proposition XI.32: Parallelepipeds of equal height are as their bases]
\label{prop:XI.32}
Parallelepipedal solids which are of the same height are to one
another as their bases.
\end{claim}
\begin{evidence}[Proof of XI.32]
\label{ev:XI.32}
By XI.25 applied with one shared dimension fixed.
\dependson{XI.32}{XI.25}
\dependson{XI.32}{XI.31}
\end{evidence}

\begin{claim}[Proposition XI.33: Similar parallelepipeds are in the triplicate ratio of corresponding edges]
\label{prop:XI.33}
Similar parallelepipedal solids are to one another in the triplicate
ratio of their corresponding sides.
\end{claim}
\begin{evidence}[Proof of XI.33]
\label{ev:XI.33}
The 3D analogue of VI.20: scaling each of the three edges by ratio
$k$ multiplies the volume by $k^3$.  Apply VI.20 to two faces and
XI.32 to extrude.
\dependson{XI.33}{VI.20}
\dependson{XI.33}{XI.32}
\dependson{XI.33}{def:V.10}
\end{evidence}

\begin{claim}[Proposition XI.34: Equal parallelepipeds have reciprocally proportional edges]
\label{prop:XI.34}
In equal parallelepipedal solids the bases are reciprocally
proportional to the heights; and those parallelepipedal solids in
which the bases are reciprocally proportional to the heights are
equal.
\end{claim}
\begin{evidence}[Proof of XI.34]
\label{ev:XI.34}
3D analogue of VI.14: by XI.32 the ratio of volumes is the
compounded ratio of bases and heights; equal volumes force the
compounded ratio to be unity, i.e.\ reciprocal proportion.
\dependson{XI.34}{VI.14}
\dependson{XI.34}{XI.32}
\dependson{XI.34}{XI.33}
\end{evidence}

\begin{claim}[Proposition XI.35: Planes equally inclined to a base have equal perpendiculars]
\label{prop:XI.35}
If there be two equal plane angles, and on their vertices there be
set up elevated straight lines containing equal angles with the
original straight lines respectively, if on the elevated straight
lines points be taken at random and perpendiculars be drawn from them
to the planes in which the original angles are, and if from the
points so arising in the planes straight lines be joined to the
vertices of the original angles, they will contain, with the elevated
straight lines, equal angles.
\end{claim}
\begin{evidence}[Proof of XI.35]
\label{ev:XI.35}
By I.4 (SAS) applied to the right triangles in each plane: equal
oblique segments and equal perpendiculars produce equal angles at
the foot.
\dependson{XI.35}{I.4}
\dependson{XI.35}{XI.10}
\dependson{XI.35}{XI.11}
\end{evidence}

\begin{claim}[Proposition XI.36: Parallelepiped on three proportionals is equal to a cube whose side is the mean]
\label{prop:XI.36}
If three straight lines be proportional, the parallelepipedal solid
formed out of the three is equal to the parallelepipedal solid on the
mean which is equilateral, but equiangular with the aforesaid solid.
\end{claim}
\begin{evidence}[Proof of XI.36]
\label{ev:XI.36}
For $a : b = b : c$ (with $b$ the mean), $abc = b^3$ and the
parallelepipeds on $\{a, b, c\}$ versus $\{b, b, b\}$ have equal
volumes by XI.34 / XI.33.
\dependson{XI.36}{VI.17}
\dependson{XI.36}{XI.32}
\dependson{XI.36}{XI.34}
\end{evidence}

\begin{claim}[Proposition XI.37: Similar parallelepipeds in proportion]
\label{prop:XI.37}
If four straight lines be proportional, the similar and similarly
described parallelepipedal solids upon them will also be
proportional; and if the similar and similarly described
parallelepipedal solids upon them be proportional, the straight lines
will themselves also be proportional.
\end{claim}
\begin{evidence}[Proof of XI.37]
\label{ev:XI.37}
3D analogue of VI.22: similar parallelepipeds are in the triplicate
ratio of edges; equality of triplicate ratios is equivalent to
equality of edge ratios.
\dependson{XI.37}{VI.22}
\dependson{XI.37}{XI.33}
\end{evidence}

\begin{claim}[Proposition XI.38: Joining diagonals of opposite faces in a cube]
\label{prop:XI.38}
If the sides of the opposite planes of a cube be bisected, and planes
be carried through the points of section, the common section of the
planes and the diameter of the cube bisect one another.
\end{claim}
\begin{evidence}[Proof of XI.38]
\label{ev:XI.38}
By symmetry of the cube under the half-turns about the centre; the
diagonal and the median plane both pass through the centre.
\dependson{XI.38}{I.10}
\dependson{XI.38}{XI.3}
\dependson{XI.38}{XI.24}
\end{evidence}

\begin{claim}[Proposition XI.39: Prisms on equal triangular bases with same height are equal]
\label{prop:XI.39}
If there be two prisms of equal height, and one have a parallelogram
as base and the other a triangle, and if the parallelogram be double
of the triangle, the prisms will be equal.
\end{claim}
\begin{evidence}[Proof of XI.39]
\label{ev:XI.39}
A triangular prism is half a parallelogram prism on the same height
(by I.34 applied to the cross-sections); the area-doubling condition
makes the two prisms have equal volume.
\dependson{XI.39}{I.34}
\dependson{XI.39}{XI.28}
\dependson{XI.39}{XI.32}
\end{evidence}
